\begin{textsample}{NEG Dim 9 – global\_north\_2025\_01 – Score -9.00 – global\_north\_2025\_01/120256432.txt}  \label{ex:f9_neg_026}
Three hostages have been freed in the first phase of the ceasefire agreement between Hamas and Israel.
The hostages, all women, were released into Red Cross custody in the Gaza Strip on Sunday and were being transferred to Israeli forces, the Israeli military said.
About 100 hostages, living and dead, are thought still to be held captive in Gaza, most of them taken in the deadly Hamas-led attack on Israel on Oct. 7, 2023.
Thirty-three of them will be released during an initial six-week phase of the ceasefire, including \textbf{female} soldiers and civilians, children, men older than 50, and sick and wounded people, according to the agreement.
Advertisement : Here is what we know about the three hostages released Sunday : Romi Gonen Gonen was 23 when she was captured as she was trying to leave the Nova \textbf{music} \textbf{festival} in southern Israel when Hamas attacked.
She was speaking at the time to her mother, Meirav Gonen, who she said had been shot and was @ @ @ @ @ @ @ @ @ @ recording of her last phone call with her \textbf{daughter}.
She told Israeli news media that Romi was a strong and happy person who often went to raves.
In the early weeks of the war, her mother expressed concern that Israeli military operations in Gaza could endanger the hostages.
Advertisement : Romi Gonen ' s older sister, Yarden, told The New York Times in February that she regularly went to a plaza in Tel Aviv, Israel, where families of hostages have held \textbf{vigils}.
Damari, 27 at the time she was captured, is the only hostage with British citizenship still being held in January.
She was taken from her home in \textbf{Kibbutz} Kfar Azza in southern Israel and was seen by a neighbor in her own car, driven by a militant, heading toward Gaza.
Damari was raised in Israel but traveled to Britain often, according to her mother, British-born Mandy Damari, who was in Israel in December to speak with officials and the news media and to plead for a @ @ @ @ @ @ @ @ @ @ been shot and that she feared for her life, telling the BBC that she had welcomed the threats from President-elect Donald Trump that there would be"all hell to pay"if no deal was reached by his inauguration.
Last January, a hostage who had been released from Gaza, Dafna Elyakim, told Israeli news media that she and her younger sister had been taken into Hamas ' underground \textbf{tunnels}, where they met other \textbf{female} hostages including Emily Damari.
Advertisement : On the eve of the anniversary of the Oct. 7 attacks, Mandy Damari spoke at an event in Hyde Park in London, where she described her \textbf{daughter} as a \textbf{soccer} fan who enjoyed a drink and had"the classic British sense of humor, with a dash of Israeli chutzpah thrown in for good measure."On Sunday, Damari thanked"everyone who never stopped fighting for Emily throughout this horrendous ordeal."But, she said in a statement,"for too many other families the impossible wait continues. @ @ @ @ @ @ @ @ @ @ when she was captured from her home in \textbf{Kibbutz} Kfar Azza, is a veterinary nurse with Romanian and Israeli citizenship.
According to Israeli news media, she was in touch with her family on the \textbf{kibbutz} when the militants attacked, telling her \textbf{parents} that they had smashed her windows and shot into her room."They ' ve arrived, they have me,"she said in a subsequent voice message sent to friends.
Last January, Hamas released a video clip of Steinbrecher and two other captives, Daniella Gilboa and Karina Ariev, in which they pleaded for their release.
In March, on her 31st birthday, the Jewish News Syndicate published an interview with her mother, Simona Steinbrecher, who said that she had looked pale and thin in the video.
She said she was concerned that her \textbf{daughter} was not getting the daily medication she needed, though she did not specify what that was.

% matched lemmas: daughter, female, festival, kibbutz, music, parent, soccer, tunnel, vigil
\end{textsample}
